\subsection{Construction pas a pas de la chaine DevSecOps de SecureRAG Hub}
\label{sec:devsecops-end-to-end}

Cette section retrace, de maniere chronologique, la construction complete de la chaine DevSecOps du projet \textit{SecureRAG Hub}. L'objectif n'est pas seulement de presenter une succession d'outils, mais de montrer comment une plateforme de chatbots metiers securises a ete progressivement industrialisee, securisee, conteneurisee, deployee localement et documentee de facon reproductible.

Le fil directeur retenu pour cette construction a ete le suivant :

\begin{enumerate}
  \item etablir un socle de depot et d'automatisation coherent ;
  \item integrer les controles de qualite et de securite des les premieres iterations ;
  \item normaliser la production des artefacts conteneurises ;
  \item deployer la plateforme sur un cluster Kubernetes local reproductible ;
  \item ajouter une chaine de confiance de type \textit{build, SBOM, sign, verify, promote, deploy} ;
  \item formaliser la CI/CD autour de Jenkins ;
  \item produire des preuves reutilisables pour la soutenance et le memoire.
\end{enumerate}

\subsubsection{Point de depart et objectifs du chantier DevSecOps}

Le projet \textit{SecureRAG Hub} repose sur une architecture microservices et sur des composants heterogenes : microservices applicatifs, orchestrateur de requetes, moteur RAG, couche securite, portail Laravel, pipeline de livraison et cluster Kubernetes local. Dans un tel contexte, la partie DevSecOps devait repondre simultanement a plusieurs contraintes :

\begin{itemize}
  \item conserver une architecture lisible et defendable dans un cadre academique ;
  \item rester executable localement sur une machine de developpement ;
  \item apporter des controles de securite concrets sur le code, les images et le cluster ;
  \item rendre possible une demonstration en soutenance sans dependre d'une infrastructure cloud externe.
\end{itemize}

Le premier principe de conception a donc ete de privilegier une pile locale, reproductible et observable, reposant sur :

\begin{itemize}
  \item un monorepo GitHub ;
  \item des scripts Bash versionnes ;
  \item Docker pour la construction des images ;
  \item \texttt{kind} pour le cluster local ;
  \item Kustomize pour la structuration des manifests ;
  \item Jenkins comme reference CI/CD finale ;
  \item Syft, Cosign, Semgrep, Gitleaks et Trivy pour les controles de securite.
\end{itemize}

\subsubsection{Etape 1 : structuration initiale du depot}

La premiere etape a consiste a structurer le depot afin de separer clairement les responsabilites :

\begin{itemize}
  \item \texttt{services/} pour les microservices applicatifs ;
  \item \texttt{infra/} pour les manifests Kubernetes, \texttt{kind} et Jenkins ;
  \item \texttt{scripts/} pour les operations CI, release, deploy, validation et secrets ;
  \item \texttt{security/} pour les regles et conventions de securite ;
  \item \texttt{docs/} pour l'architecture, les runbooks et les sections memoire ;
  \item \texttt{platform/} pour la plateforme Laravel et ses composants associes.
\end{itemize}

Cette organisation a constitue le premier jalon de la maturite DevSecOps, car elle a permis d'eviter un depot monolithique mal compartimente et d'associer rapidement chaque brique technique a un emplacement stable et documente.

\subsubsection{Etape 2 : mise en place des controles CI de base}

Une fois le depot structure, le chantier s'est concentre sur les verifications de qualite et de securite qui devaient etre executees regulierement.

\paragraph{Tests et couverture.}
Des scripts dedies ont ete ajoutes pour executer les tests automatisees et mesurer la couverture, avec un objectif pedagogique simple : disposer d'un premier niveau de confiance avant toute industrialisation plus lourde.

\paragraph{SAST et detection de secrets.}
Le depot a ensuite integre :

\begin{itemize}
  \item \texttt{Semgrep} pour la detection statique de patterns de code a risque ;
  \item \texttt{Gitleaks} pour la recherche de secrets ou de traces de credentials dans le depot ;
  \item \texttt{Trivy} pour les scans orientes fichiers et images.
\end{itemize}

Cette etape a ete importante d'un point de vue methodologique : plutot que d'ajouter la securite uniquement au moment du deploiement, le projet a choisi de l'integrer des la phase de verification continue.

\subsubsection{Etape 3 : conteneurisation des microservices}

La phase suivante a consiste a normaliser la construction des artefacts deployables. Chaque microservice a ete accompagne d'un \texttt{Dockerfile}, avec l'objectif de produire des images coherentes, reutilisables localement et compatibles avec une execution sur Kubernetes.

Les services metiers principaux concernes sont :

\begin{itemize}
  \item \texttt{auth-users}
  \item \texttt{chatbot-manager}
  \item \texttt{conversation-service}
  \item \texttt{audit-security-service}
  \item \texttt{portal-web}
\end{itemize}

Les anciens dossiers Python/RAG sous \texttt{services/} sont conserves comme historique technique, mais exclus de la chaine officielle tant que leurs sources applicatives ne sont pas restaurees.

\subsubsection{Etape 4 : premier socle Kubernetes avec kind et Kustomize}

Le premier vrai saut d'industrialisation a ete l'introduction d'un cluster local \texttt{kind}. Ce choix a repondu a trois objectifs :

\begin{itemize}
  \item reproduire un environnement Kubernetes realiste ;
  \item permettre un deploiement integralement local ;
  \item maitriser le cycle complet \textit{build, push, deploy, validate} sur une machine unique.
\end{itemize}

Le deploiement a ete structure autour de Kustomize, avec une separation en \texttt{base/overlay}. La base contient les objets communs et les manifests de services ; les overlays permettent d'adapter le comportement selon le contexte, en particulier \texttt{dev} et \texttt{demo}.

Les premiers composants de la couche Kubernetes ont ete :

\begin{itemize}
  \item le namespace \texttt{securerag-hub} ;
  \item les \textit{ServiceAccount} dedies ;
  \item les \textit{Deployment}, \textit{Service} et \textit{StatefulSet} ;
  \item les \textit{ConfigMap} et \textit{Secret} ;
  \item les premiers \textit{NetworkPolicy} ;
  \item les probes \texttt{readiness}, \texttt{liveness} et ensuite \texttt{startup}.
\end{itemize}

\subsubsection{Etape 5 : durcissement minimal des workloads}

Une fois le socle \texttt{kind + Kustomize} fonctionnel, l'effort a porte sur le durcissement des conteneurs et la reduction de la surface d'attaque.

Les choix principaux ont ete les suivants :

\begin{itemize}
  \item \texttt{runAsNonRoot} lorsque l'image le permet ;
  \item \texttt{allowPrivilegeEscalation=false} ;
  \item \texttt{readOnlyRootFilesystem=true} lorsque possible ;
  \item \texttt{capabilities.drop: [ALL]} ;
  \item \texttt{seccompProfile: RuntimeDefault} ;
  \item \textit{ServiceAccount} dedie par composant ;
  \item \textit{default deny} reseau puis ouvertures minimales ciblees.
\end{itemize}

Cette etape a permis de transformer un simple deploiement de conteneurs en un premier environnement Kubernetes securise et defendable techniquement.

\subsubsection{Etape 6 : gestion du registre local et execution locale complete}

Pour eviter toute dependance forte a un registre externe, un registre local \texttt{localhost:5001} a ete introduit. La logique retenue a ete la suivante :

\begin{enumerate}
  \item construire localement les images OCI ;
  \item les pousser dans un registre local ;
  \item configurer \texttt{kind} pour les consommer ;
  \item deployer ensuite les manifests Kubernetes.
\end{enumerate}

Cette architecture locale a un interet academique fort : elle permet de demontrer la chaine complete d'industrialisation sans infrastructure cloud, tout en restant proche d'un fonctionnement professionnel.

\subsubsection{Etape 7 : scripts de deploiement et de validation}

Afin de fiabiliser l'usage quotidien du projet, des scripts operationnels ont ete ajoutes :

\begin{itemize}
  \item creation du cluster et du registre ;
  \item build des images ;
  \item chargement des images dans \texttt{kind} ;
  \item deploiement Kustomize ;
  \item smoke tests ;
  \item validations fonctionnelles, securite et RAG ;
  \item generation de rapports post-deploiement.
\end{itemize}

Cette etape a marque le passage d'un ensemble de manifests a une veritable procedure de livraison locale.

\subsubsection{Etape 8 : premiere campagne de validation post-deploiement}

Une fois le runtime local deployable, plusieurs scripts de validation ont ete introduits pour couvrir :

\begin{itemize}
  \item la disponibilite des services ;
  \item les endpoints internes \texttt{/healthz} ;
  \item l'acces externe par \texttt{NodePort} ;
  \item les composants RAG ;
  \item la couche \texttt{security-auditor} ;
  \item la production d'un rapport synthetique \texttt{PASS / FAIL / SKIP}.
\end{itemize}

Le choix du triplet \texttt{PASS}, \texttt{FAIL}, \texttt{SKIP} a ete deliberement retenu pour s'adapter a un projet en construction. Il permet de distinguer une anomalie technique reelle d'un endpoint metier non encore implemente.

\subsubsection{Etape 9 : construction de la supply chain securisee}

Le chantier DevSecOps a ensuite ete etendu a la \textit{supply chain security}. L'objectif n'etait plus seulement de deployer des images, mais de garantir la confiance associee a leur production.

Trois briques principales ont ete introduites :

\begin{itemize}
  \item generation de SBOM CycloneDX avec \texttt{Syft} ;
  \item signature d'images OCI avec \texttt{Cosign} ;
  \item verification des signatures avant promotion ou deploiement.
\end{itemize}

Le projet a egalement adopte une convention de nommage d'images stable :

\begin{verbatim}
REGISTRY_HOST/IMAGE_PREFIX-<service>:IMAGE_TAG
\end{verbatim}

Cette convention a rendu possible un fonctionnement coherent des scripts de build, des scripts de release et des overlays Kubernetes.

\subsubsection{Etape 10 : normalisation de la release et promotion sans reconstruction}

Une fois les scripts SBOM, signature et verification disponibles, la question centrale est devenue : comment eviter un \textit{rebuild} en phase de deploiement ?

La reponse apportee par le projet a ete de formaliser la boucle de confiance suivante :

\begin{center}
\texttt{build -> sbom -> sign -> verify -> promote -> deploy}
\end{center}

Concretement :

\begin{itemize}
  \item la phase de release produit un artefact image ;
  \item le SBOM est genere pour cet artefact ;
  \item l'image est signee ;
  \item la signature est verifiee ;
  \item un script de promotion cree ensuite le tag cible sans reconstruire ;
  \item le deploiement s'appuie sur le tag ou digest promu.
\end{itemize}

Ce choix constitue une bonne pratique professionnelle majeure, car il elimine l'ambiguite entre image testee et image deployee.

\subsubsection{Etape 11 : transition de GitHub Actions vers Jenkins}

Le projet disposait initialement de workflows GitHub Actions. Toutefois, pour disposer d'une chaine CI/CD plus explicite, plus demonstrable localement et plus conforme a une soutenance technique, Jenkins a ete retenu comme reference finale.

Cette transition a eu plusieurs consequences :

\begin{itemize}
  \item separation explicite entre \texttt{Jenkinsfile} pour la CI et \texttt{Jenkinsfile.cd} pour la CD ;
  \item maintien des scripts shell comme coeur de la logique d'execution ;
  \item reduction de la logique enfouie dans l'orchestrateur CI/CD ;
  \item possibilite de rejouer les memes commandes hors Jenkins.
\end{itemize}

Cette decision a permis de conserver une approche \textit{pipeline as code} tout en rendant la demonstration plus lisible et plus autonome.

\subsubsection{Etape 12 : industrialisation de Jenkins en local}

L'etape suivante a consiste a faire de Jenkins un composant reellement exploitable, et non un simple \textit{Jenkinsfile} theorique. Une stack locale Jenkins a donc ete ajoutee avec :

\begin{itemize}
  \item un \texttt{docker-compose} dedie ;
  \item une liste de plugins maitrisee ;
  \item une configuration \textit{Configuration as Code} ;
  \item des jobs DSL versionnes ;
  \item une strategie de credentials locaux ;
  \item des scripts de bootstrap et de collecte de preuves.
\end{itemize}

L'objectif etait double :

\begin{itemize}
  \item prouver que Jenkins peut etre lance localement et seeded automatiquement ;
  \item rendre reproductible la creation des jobs \texttt{securerag-hub-ci} et \texttt{securerag-hub-cd}.
\end{itemize}

Cette etape a completement change le niveau de maturite du projet : la CI/CD n'est plus seulement definie dans le code, elle devient une infrastructure locale executable et versionnee.

\subsubsection{Etape 13 : gestion des secrets}

Le traitement des secrets a ete formalise a trois niveaux.

\paragraph{Niveau local.}
Des fichiers exemple et des scripts de bootstrap ont ete ajoutes pour generer les secrets necessaires au developpement local, sans les committer dans le depot.

\paragraph{Niveau Jenkins.}
Les credentials Jenkins ont ete normalises, en particulier pour :

\begin{itemize}
  \item \texttt{cosign-private-key}
  \item \texttt{cosign-public-key}
  \item \texttt{cosign-password}
\end{itemize}

\paragraph{Niveau Kubernetes.}
Les manifests applicatifs s'appuient sur des \textit{Secret} references via \texttt{secretRef}, et des scripts permettent de creer les secrets de developpement du namespace \texttt{securerag-hub}.

Le point important ici est la separation claire entre :

\begin{itemize}
  \item exemples et gabarits versionnes ;
  \item credentials d'exploitation non versionnes ;
  \item secrets injectes au runtime.
\end{itemize}

\subsubsection{Etape 14 : stabilisation du runtime local et gestion du cas Ollama}

Le composant \texttt{Ollama} a constitue le point le plus sensible du runtime local. Son image lourde et son comportement plus lent au premier demarrage ont cree un risque pour une demonstration fluide.

Plutot que de masquer ce probleme, le projet a adopte une strategie plus robuste :

\begin{itemize}
  \item conserver un mode \texttt{dev} avec le composant reel ;
  \item ajouter un overlay \texttt{demo} permettant un fonctionnement fallback ;
  \item documenter ce mode comme une strategie de demonstration, et non comme une substitution de production.
\end{itemize}

Cette approche est importante d'un point de vue professionnel : elle montre qu'un environnement de demonstration n'a pas besoin d'etre fragile pour etre techniquement honnete, a condition que les ecarts soient explicites et documentes.

\subsubsection{Etape 15 : integration du portail Laravel dans le runtime Kubernetes}

Initialement, le portail Laravel etait surtout demarrable localement. Dans un deuxieme temps, il a ete integre comme workload Kubernetes a part entiere, avec :

\begin{itemize}
  \item \textit{Deployment}
  \item \textit{Service}
  \item \textit{ServiceAccount}
  \item \textit{NetworkPolicy}
  \item exposition via \texttt{NodePort}
\end{itemize}

Cette integration a renforce la coherence globale du projet : la plateforme web ne reste plus en marge de la chaine DevSecOps, elle fait partie du meme socle de build, release, deploiement et validation.

\subsubsection{Etape 16 : documentation d'exploitation et de soutenance}

En parallele des composants techniques, un ensemble de documents a ete construit afin de rendre le projet transmissible et defendable :

\begin{itemize}
  \item architecture cible DevSecOps ;
  \item runbook Jenkins ;
  \item runbook \texttt{kind} ;
  \item runbook de promotion release ;
  \item runbook de troubleshooting ;
  \item strategie secrets ;
  \item matrice de controles ;
  \item rapports LaTeX pour le memoire.
\end{itemize}

Ce travail documentaire est central dans un projet de niveau Master DSIR. Il demontre que la valeur du chantier DevSecOps ne repose pas uniquement sur l'automatisation, mais aussi sur la capacite a expliquer, transmettre, auditer et reproduire la solution.

\subsubsection{Etape 17 : corrections appliquees et fiabilisation finale}

La phase finale du chantier n'a pas seulement consiste a ajouter de nouveaux composants. Elle a aussi permis de corriger plusieurs problemes reels observes lors des executions locales, afin d'ameliorer la demonstrabilite du projet sans masquer les limites de l'environnement.

\paragraph{Correction de l'environnement Python local.}
Lors des premiers essais de la commande \texttt{make test}, les tests FastAPI echouaient non pas a cause du code applicatif, mais a cause d'un probleme d'architecture entre l'interpreteur Python local et certaines extensions binaires, notamment \texttt{pydantic\_core} et \texttt{coverage}. La machine de test etant en \texttt{arm64}, alors que certains paquets avaient ete installes pour \texttt{x86\_64}, la collecte des tests s'interrompait avant execution.

La correction retenue a ete simple et professionnelle :

\begin{itemize}
  \item creation d'un environnement virtuel dedie au depot ;
  \item reinstallation propre des dependances dans cet environnement ;
  \item abandon de l'utilisation du Python global comme source de verite pour les tests.
\end{itemize}

Cette correction a permis de retablir une campagne de tests locale cohérente, avec \texttt{12/12} tests au statut \texttt{PASS} et un taux de couverture mesure a \texttt{100\%}.

\paragraph{Correction de la chaine de signature locale.}
Un second ecart a ete observe lors des essais de verification et de promotion d'images : l'outil \texttt{Cosign} etait absent de l'environnement local. Les scripts de release fonctionnaient donc correctement d'un point de vue logique, mais leur execution etait bloquee par un pre-requis non satisfait.

La correction appliquee a consiste a :

\begin{itemize}
  \item installer \texttt{Cosign} localement ;
  \item utiliser les scripts de bootstrap du depot pour generer une paire de cles locale ;
  \item normaliser l'usage des variables \texttt{COSIGN\_KEY}, \texttt{COSIGN\_PUBLIC\_KEY} et \texttt{COSIGN\_PASSWORD}.
\end{itemize}

Cette etape est importante dans le memoire car elle montre que la securisation de la release ne depend pas seulement des scripts, mais aussi d'une hygiene explicite de l'environnement d'execution.

\paragraph{Correction de l'interaction entre Kyverno et les pods de validation.}
L'activation reelle de Kyverno a fait apparaitre un troisieme point de vigilance. Les scripts de validation post-deploiement demarraient un pod ephemere en utilisant une image \texttt{securerag-hub-*} du registre local \texttt{localhost:5001}. Cette image entrait donc dans le perimetre de la policy de verification Cosign. Or, au moment du test, le registre local n'etait pas joignable par Kyverno via \texttt{localhost} depuis le contexte du cluster, ce qui provoquait un refus d'admission.

La correction retenue n'a pas ete de desactiver la policy, mais de clarifier le protocole de validation :

\begin{itemize}
  \item conserver la policy Kyverno en mode \texttt{Audit} par defaut ;
  \item reserver le mode \texttt{Enforce} aux environnements ou la chaine de signature est effectivement stabilisee ;
  \item utiliser une image publique de validation, par exemple \texttt{python:3.12-slim}, pour les pods ephemeres lorsque l'objectif est de tester les endpoints internes sans faire dependre ces tests d'un registre local signe.
\end{itemize}

Cette correction est pedagogiquement utile car elle distingue clairement deux objectifs differents :

\begin{itemize}
  \item verifier la securite de la chaine d'images applicatives ;
  \item verifier la sante du cluster et des endpoints internes.
\end{itemize}

\paragraph{Activation effective des addons de demonstration.}
La fin du chantier a egalement permis de transformer deux briques jusque-la seulement preparees en composants reellement activables sur le cluster :

\begin{itemize}
  \item \texttt{metrics-server}, afin de rendre les \textit{Horizontal Pod Autoscaler} effectivement observables ;
  \item \texttt{Kyverno}, afin de rendre les policies d'admission applicables au cluster local.
\end{itemize}

Une fois \texttt{metrics-server} installe, le \textit{HPA} de \texttt{portal-web} doit etre observable avec des valeurs CPU reelles. De meme, une fois \texttt{Kyverno} installe, les \textit{ClusterPolicy} du projet doivent etre detectees a l'etat \texttt{Ready}. Sans les sorties \texttt{kubectl top}, \texttt{kubectl get hpa}, \texttt{kubectl get clusterpolicy} et \texttt{kubectl get policyreport -A}, ces controles restent \texttt{DEPENDANT\_DE\_L\_ENVIRONNEMENT}.

\paragraph{Consolidation honnete de la preuve finale.}
Enfin, la campagne de reference unique a ete renforcee pour produire un mode \texttt{dry-run}. Ce choix n'a pas pour but de simuler un succes fictif, mais de fournir un support demonstratif honnete lorsque certains pre-requis ne sont pas encore reunis, par exemple l'absence d'images deja signees dans le registre local au moment de la soutenance.

La preuve finale du projet distingue donc desormais explicitement :

\begin{itemize}
  \item ce qui a ete reellement execute et observe ;
  \item ce qui est pret a etre rejoue ;
  \item ce qui reste conditionne par l'etat local de Docker, du registre, de \texttt{kind} ou des cles Cosign.
\end{itemize}

\subsubsection{Etat final atteint}

Au terme du chantier, les elements suivants sont consideres comme finalises dans le depot :

\begin{itemize}
  \item Jenkins local ;
  \item jobs CI/CD ;
  \item scripts de release ;
  \item scripts de deploiement verifie ;
  \item documentation de type runbook ;
  \item strategie de secrets ;
  \item overlay \texttt{demo} ;
  \item integration de \texttt{portal-web} dans Kubernetes.
\end{itemize}

Autrement dit, la chaine DevSecOps de \textit{SecureRAG Hub} couvre desormais :

\begin{itemize}
  \item la verification continue du code ;
  \item la production d'artefacts OCI ;
  \item la generation d'un SBOM ;
  \item la signature et la verification des images ;
  \item la promotion d'artefacts sans reconstruction ;
  \item le deploiement sur un cluster local \texttt{kind} ;
  \item la validation post-deploiement ;
  \item la collecte de preuves techniques pour le memoire et la soutenance.
\end{itemize}

\subsubsection{Analyse critique de la demarche}

La construction de la chaine DevSecOps de \textit{SecureRAG Hub} permet de tirer plusieurs enseignements.

\paragraph{Premier enseignement : la reproductibilite est aussi importante que l'automatisation.}
Une pile locale reproductible, basee sur scripts shell, Docker, \texttt{kind} et Jenkins, s'est revelee plus defendable pedagogiquement qu'une solution tres distribuee dependant d'equipements externes.

\paragraph{Deuxieme enseignement : la securite doit etre repartie tout au long du cycle de vie.}
Le projet n'a pas concentre la securite dans un unique scan ou dans une unique etape de deploiement. Au contraire, il a reparti les controles sur le code, les images, la signature, les manifests, le runtime et la validation finale.

\paragraph{Troisieme enseignement : la documentation fait partie de la valeur DevSecOps.}
Les runbooks, les rapports et les annexes ne sont pas un simple complement de forme. Ils constituent la preuve que la solution peut etre transmise, reprise et evaluee.

\subsubsection{Conclusion generale}

La partie DevSecOps de \textit{SecureRAG Hub} a ete construite de facon progressive, en partant d'un depot structure et de scripts simples, puis en evoluant vers une chaine complete integree a Jenkins, a Kubernetes local et a une logique de \textit{supply chain security}. Le projet atteint ainsi un niveau de maturite satisfaisant pour un PFE Master DSIR : il est techniquement credible, localement executable, securise de maniere concrete et suffisamment documente pour etre soutenu de facon rigoureuse.

Cette progression pas a pas constitue un resultat en soi. Elle montre que l'industrialisation securisee d'une plateforme applicative ne se limite pas au choix de quelques outils ; elle releve d'une demarche d'architecture, de normalisation, de verification et de preuve, menee sur l'ensemble du cycle de vie logiciel.
